Symbolic regression aims to recover closed-form expressions
from numerical data, but in differentiable symbolic regression
the recovered expression depends not only on the grammar but
also on the fixed architecture through which variables are
routed during training.  This is relevant to signal-processing
settings in which closed-form models and interpretable nonlinear
structure are useful.  This architecture-specific effect has
rarely been isolated directly, because existing comparisons often
vary architecture together with operator family, grammar, or
search procedure.
Three depth-3 architectures are compared across twenty-four
operator--shape--leaf combinations, holding operator family,
grammar, and training protocol fixed as far as possible while
varying the variable-routing architecture.  Recovery changes from
$0/64$ to $64/64$ trials on the same target under an
architecture-plus-native-training-protocol comparison.  The best
architecture on one target is the worst on another, and trees
with two equal-depth subtrees fail in every configuration tested
($0/3{,}776$).  As a proof-of-concept mitigation, a small
architecture set is trained and the hardened expression with the
lowest held-out RMSE is selected.  On the jointly-run subset,
this improves recovery from $34.4\%$ for the only architecture
present in all three configurations to $50.1\%$.  On a
Shockley diode target, the validation selector recovers cases
missed by that baseline architecture, which by itself recovers
$0/32$ seeds.  Since the jointly-run subset contains only three
configurations, the selector result is evidence that
validation-based architecture selection is promising, not a
complete benchmark.  These results support treating architecture
as a measurable design variable that should be reported,
stress-tested, and selected using held-out validation rather
than fixed a priori.
